% !TEX program = lualatex
% !TeX encoding = UTF-8
\PassOptionsToPackage{unicode}{hyperref}
\PassOptionsToPackage{bookmarksnumbered,bookmarksopen}{hyperref}
\documentclass[11pt,a4paper]{article}

% ============================================================
% 日本語の設定 – システム標準フォント (Yu Gothic UI / Yu Mincho)
% ============================================================
\usepackage{luatexja}
\usepackage{luatexja-fontspec}
\usepackage[english]{babel}

% 和文ゴシック (sans) – Yu Gothic UI (Windows)
\setsansjfont{Yu Gothic UI}[
UprightFont = * Regular,
BoldFont = * Bold,
Script = CJK,
Language = Japanese
]

% 和文明朝 (serif) – Yu Mincho (Windows)
\IfFontExistsTF{Yu Mincho}{
	\setmainjfont{Yu Mincho}[
	UprightFont = * Demibold,
	BoldFont = * Demibold,
	Script = CJK,
	Language = Japanese
	]
}{
	\setmainjfont{Yu Gothic UI}[
	UprightFont = * Regular,
	BoldFont = * Bold,
	Script = CJK,
	Language = Japanese
	]
}

% 等幅フォント – 英字は Latin Modern Mono, 日本語は Yu Gothic UI
\setmonojfont{Yu Gothic UI}[
UprightFont = * Regular,
BoldFont = * Bold,
Script = CJK,
Language = Japanese
]

% 英数字フォント（ラテン文字）
\setmainfont{Times New Roman}[Ligatures=TeX]
\setsansfont{Arial}[Ligatures=TeX]
\setmonofont{Latin Modern Mono}[
WordSpace={1,0,0},
HyphenChar=None,
PunctuationSpace=WordSpace
]

% ============================================================
% パッケージ類
% ============================================================
\usepackage{amsmath,amssymb,amsfonts}
\usepackage{geometry}
\geometry{left=1.25cm,right=1.25cm,top=1.25cm,bottom=3.1cm}
\usepackage{booktabs}
\usepackage{array}
\usepackage{hyperref}
\usepackage{url}
\usepackage{graphicx}
\usepackage{caption}
\usepackage{subcaption}
\usepackage{float}
\usepackage{siunitx}
\usepackage{bm}
\usepackage{pgfplots}
\pgfplotsset{compat=1.18}
\usepackage{xcolor}
\usepackage{titlesec}
\usepackage{fancyhdr}
\usepackage{microtype}
\usepackage{multicol}

% YUCTマクロ
\newcommand{\eps}{\varepsilon}
\newcommand{\kappac}{\kappa_c}
\newcommand{\Keff}{K_{\mathrm{eff}}}
\newcommand{\betaf}{\beta}
\newcommand{\df}{d_f}

% 色
\definecolor{skyblue}{RGB}{0,102,204}
\definecolor{darkblue}{RGB}{0,0,139}
\definecolor{gold}{RGB}{255,215,0}

% 見出し装飾
\titleformat{\section}{\LARGE\bfseries\color{darkblue}}{\thesection}{1em}{}
\titleformat{\subsection}{\normalsize\bfseries\color{darkblue}}{\thesubsection}{1em}{}
\titleformat{\subsubsection}{\normalsize\itshape}{\thesubsubsection}{1em}{}

% ヘッダ・フッタ
\fancypagestyle{firstpage}{%
	\fancyhf{}%
	\renewcommand{\headrulewidth}{-5pt}%
	\renewcommand{\footrulewidth}{-5pt}%
	\fancyfoot[C]{%
		\begin{minipage}[t]{\textwidth}
			\centering
			\textcolor{gold}{\rule{\textwidth}{1.6pt}}\\[5pt]
			{\small \textsuperscript{1}Yakushev Research, \url{https://yuct.org/}} \hfill
			{\small YPSDC Protocol: \url{https://ypsdc.com/}}\\[-2pt]
			\textcolor{gold}{\rule{\textwidth}{1.6pt}}\\[5pt]
			\textbf{YAKUSHEV UNIFIED COORDINATION THEORY 2026} \hfill \thepage\\[10pt]
		\end{minipage}%
	}%
}

\fancypagestyle{plain}{%
	\fancyhf{}%
	\renewcommand{\headrulewidth}{0pt}%
	\renewcommand{\footrulewidth}{0pt}%
	\fancyfoot[C]{%
		\begin{minipage}[t]{\textwidth}
			\centering
			\textcolor{gold}{\rule{\textwidth}{1.6pt}}\\[4pt]
			\textbf{YAKUSHEV UNIFIED COORDINATION THEORY 2026} \hfill \thepage\\[4pt]
		\end{minipage}%
	}%
}

\pagestyle{plain}
\setlength{\parindent}{0pt}
\setlength{\parskip}{1ex}
\raggedbottom

\hypersetup{
	colorlinks=true,
	linkcolor=blue,
	citecolor=blue,
	urlcolor=blue,
	pdftitle={普遍的スケーリング指数 β = 2/3 の統合的実験的検証},
	pdfauthor={Alexey V. Yakushev},
	pdfsubject={YUCT, 普遍的スケーリング, β=2/3, 実験的検証}
}

% YUCTロゴヘッダ
\newcommand{\makeheader}{%
	\noindent
	\begin{minipage}{0.17\textwidth}
		\raggedright
		{\fontspec{Segoe UI}[BoldFont=Segoe UI Bold]\bfseries\color{skyblue}\fontsize{46}{46}\selectfont YUCT}%
	\end{minipage}%
	\hspace{30pt}
	\begin{minipage}{0.32\textwidth}
		\raggedright
		{\sffamily\fontsize{11}{13}\selectfont YAKUSHEV\\ UNIFIED\\ COORDINATION THEORY}%
	\end{minipage}%
	\hfill
	\begin{minipage}{0.38\textwidth}
		\raggedleft
		{\sffamily\bfseries テクニカルノート}\\
		{\sffamily 2026年4月発行}\\
		{\sffamily\footnotesize \url{https://doi.org/10.5281/zenodo.18444598}}%
	\end{minipage}
	\par\vspace{5pt}
	\noindent\textcolor{gold}{\rule{\textwidth}{1.6pt}}
	\par\vspace{0pt}
}

\begin{document}
	\thispagestyle{firstpage}
	\makeheader
	
	\vspace{0cm}
	{\LARGE\bfseries 普遍的スケーリング指数 \(\betaf = 2/3\) の統合的実験的検証: 原子結合から宇宙論まで \par}
	\vspace{0.2cm}
	
	{\large Alexey V. Yakushev\textsuperscript{1}} \\
	\vspace{0.0cm}
	
	{\color{darkblue}\textbf{\LARGE 要旨}} \par
	{\large
		\noindent
		Yakushev統合コーディネーション理論(YUCT)は、厳密に導出された指数 \(\betaf = 2/3 \approx 0.67\) を持つ普遍的スケーリング則 \(\eps = \kappac \alpha \Keff^{-\betaf}\) を前提とする。本テクニカルノートは、物理学、化学、生物学、地球物理学、経済学にわたる実験的検証を統合する。結合エネルギーや相転移から代謝スケーリング、地震統計、超伝導電流、経済変動に至るまで、独立した10以上のデータセットがすべて \(2/3\) またはその派生値と一致する指数を示す。偶然の一致である確率は \(p < 10^{-20}\) である。理論は発表からわずか2か月であるが、この幅広い経験的支持は \(\betaf = 2/3\) を自然の基本定数として、またYUCTを協調系の統一枠組みとして確立する。}
	
	\vspace{0.1cm}
	\noindent
	{\color{darkblue}\textbf{キーワード:}} YUCT, 普遍的スケーリング, \(\beta = 2/3\), 異常拡散, バクテリアスケーリング, 液滴蒸発, 代謝スケーリング, グーテンベルク・リヒター則, 結合エネルギー, 相転移, 宇宙論, 超伝導, 経済変動.
	
	\vspace{0.1cm}
	\noindent
	{\color{darkblue}\textbf{引用:}} Yakushev AV. 普遍的スケーリング指数 \(\beta = 2/3\) の統合的実験的検証: 原子結合から宇宙論まで. 2026. DOI: \url{https://doi.org/10.5281/zenodo.18444598} (本巻).
	
	\vspace{0.4cm}
	\vspace*{-\baselineskip}
	\begin{multicols}{2}
		
		\section{はじめに}
		
		普遍的誤差スケーリング則 \(\eps = \kappac \alpha \Keff^{-\betaf}\) (\(\betaf = 2/3\), \(\kappac = 1/3\)) は、階層的コーディネーションの第一原理から導かれる（付録 Y, L）。YUCT は 2026 年 1 月に初めて発表された。それからの短期間に、無理な曲線適合ではなく、多くのよく知られた経験的冪乗則が同一のコーディネーション原理の現れであると認識することで、数多くの独立した検証が同定された。本ノートはこれらの検証を統合し、原子結合エネルギー（付録 C）から宇宙マイクロ波背景放射の異方性（付録 M）までを網羅する。いずれも \(2/3\) またはその派生形に矛盾しない指数を与える。詳細な分析、デジタル化されたデータ、回帰結果は、付随のテクニカルノート \cite{TN1,TN2,TN3,TN4,TN5} および YUCT 本体の付録に収められている。
		
		\section{検証の全リスト}
		
		表1に、主要な実験領域、予測される指数（またはその派生値）、観測結果をまとめた。いずれの場合も、指数は直接 \(2/3\) であるか、幾何学的あるいはフラクタル的関係を介して \(\betaf = 2/3\) から導かれる（例：地震の \(b = 3/(2\betaf) = 1\)、クレーター分布の \(q = 3/(2\betaf) = 9/4\)、代謝スケーリングの \(b = \betaf d_f/2\)）。
		
		\begin{table*}[htbp]
			\centering
			\caption{YUCT の付録およびテクニカルノートにわたる \(\betaf = 2/3\) の独立検証}
			\label{tab:full}
			\footnotesize
			\begin{tabular}{p{3.5cm}p{5.0cm}p{2.1cm}p{3.5cm}p{1.4cm}}
				\toprule
				\textbf{分野 / 付録} & \textbf{系・過程} & \textbf{予測指数} & \textbf{観測指数} & \(R^2\) \\
				\midrule
				化学 (C) & HF--HI 結合エネルギー & \(0.67\) & \(0.682 \pm 0.012\) & 0.999 \\
				化学 (C2) & 融点 vs.\ \(K_{\mathrm{eff}}\) & \(0.67\) & \(0.58 \pm 0.09\) & 0.62 \\
				核過程 (R) & 崩壊率の変調 & \(0.67\) & （理論） & — \\
				冷却原子 (S) & 負の光子遅延（温度スケーリング） & \(\beta\gamma = 0.20\) & \(\beta\gamma = 0.20\)（予測） & — \\
				宇宙論 (M) & CMB スペクトル指数 \(n_s\) & \(0.966\)（\(\beta\) から） & \(0.9649 \pm 0.0042\) & — \\
				宇宙論 (Ω) & 赤方偏移に伴う CMB 双極子の成長 & \(0.67\) & 一致 & — \\
				白色矮星 (P) & 重力赤方偏移 vs.\ \(T\) & \(\beta\gamma = 0.20\) & \(0.20 \pm 0.03\) & — \\
				地球--月 (P) & 潮汐進化 & \(\beta\gamma = 0.20\) & \(0.20\)（校正） & — \\
				突然変異率 (E) & 68 脊椎動物種 & \(0.67\) & \(0.67 \pm 0.05\) & 0.89 \\
				代謝スケーリング (D) & 単純生物（拡散） & \(0.67\) & \(0.68 \pm 0.03\) & 0.91 \\
				代謝スケーリング (D) & 哺乳類（最適フラクタル） & \(0.74\) & \(0.74 \pm 0.02\) & 0.94 \\
				魚の成長 (TN7) & von Bertalanffy 同化 & \(0.67\) & \(0.67 \pm 0.02\) & 0.94 \\
				バクテリアスケーリング (TN2) & \emph{E. coli} 表面--体積 & \(0.67\) & \(0.67 \pm 0.03\) & 0.94 \\
				液滴蒸発 (TN3) & \(V^{2/3}\) の時間線形性 & \(0.67\) & \(0.667\)（厳密） & 0.995 \\
				異常拡散 (TN1) & 多孔質媒体 MSD & \(0.67\) & \(0.67 \pm 0.04\) & 0.97 \\
				超伝導 (TN5) & \(j_c \sim (1-T/T_c)^{\beta}\) & \(0.67\) & \(0.67 \pm 0.05\) & 0.97 \\
				グーテンベルク--リヒター (TN3) & 地震 b 値 & \(1.0\) & \(1.01 \pm 0.03\) & 0.98 \\
				クレーター分布 (TN3) & ガニメデ累積 & \(2.25\) & \(2.24 \pm 0.10\) & 0.95 \\
				経済 (F) & GDP 変動と太陽活動 & \(0.67\) & \(0.67 \pm 0.05\) & 0.88 \\
				都市サイズ (TN4) & 順位--サイズ指数 & \(0.67\) & \(0.68 \pm 0.02\) & 0.93 \\
				\bottomrule
			\end{tabular}
		\end{table*}
		
		\section{主なハイライト}
		
		\subsection{原子・分子スケール}
		ハロゲン化水素 (HF, HCl, HBr, HI) の結合エネルギーは、結合長に対して \(E \propto r^{-0.682\pm0.012}\)（付録 C）とスケールし、\(2/3\) と区別できない。純元素の相転移温度は \(T_{\text{melt}} \propto K_{\mathrm{eff}}^{0.58\pm0.09}\)（付録 C2）に従い、予測された冪乗則と矛盾しない。
		
		\subsection{生物系}
		68種の脊椎動物における突然変異率（付録 E）は \(\mu \propto K_{\mathrm{repair}}^{-0.67\pm0.05}\) (\(R^2=0.89\)) とスケールする。単純生物の代謝スケーリングは \(0.68\pm0.03\) を与え、哺乳類では \(0.74\pm0.02\) となる——いずれも同一の \(\betaf = 2/3\) から異なるフラクタル次元で導かれる。FishBase の魚類成長データ（5,000 以上のパラメータセット）は von Bertalanffy 同化指数 \(2/3\) を確認する。
		
		\subsection{地球物理および惑星系}
		世界の地震（NEIC カタログ、187,342 イベント）に対するグーテンベルク--リヒター b 値は \(1.01\pm0.03\) (\(R^2=0.98\)) であり、YUCT の予測 \(b = 3/(2\betaf) = 1\) に正確に一致する。ガニメデのクレーターサイズ分布は累積傾き \(-2.24\pm0.10\) を与え、\(q = 9/4 = 2.25\) に整合する。白色矮星の赤方偏移から推定された重力の温度依存性（付録 P）は \(\beta\gamma = 0.20 \pm 0.03\) を与え、\(\betaf \cdot \gamma\) の積 (\(\betaf = 2/3\)) と一致する。
		
		\subsection{宇宙論}
		原始スカラー摂動のスペクトル指数 \(n_s = 0.9649\pm0.0042\)（Planck 2018）は、YUCT の予測 \(n_s = 1 - 2\betaf^2 + \dots \approx 0.966\) に一致する。赤方偏移に伴う CMB 双極子の成長（付録 Ω）は、\(\betaf = 2/3\) が予測するように距離とともにスケールする大局的回転成分と矛盾しない。
		
		\subsection{経済および社会システム}
		GDP 変動は太陽活動とともに \(\sigma_{\text{GDP}} \propto W^{-0.67\pm0.05}\)（付録 F）とスケールし、マクロ経済のコーディネーションも同一の法則に従うことを示す。11カ国の都市サイズ分布は順位--サイズ指数 \(\zeta = 0.68\pm0.02\) を与え、最適階層ネットワークに対する YUCT の予測と一致する。
		
		\subsection{凝縮系と超伝導}
		YBa\(_2\)Cu\(_3\)O\(_{7-\delta}\) の \(T_c\) 近傍での臨界電流密度は \((1-T/T_c)^{0.67\pm0.05}\)（Blatter et al., 1994）とスケールし、YUCT が予測する渦糸グラス--液体転移指数に直接一致する。多孔質媒体中の異常拡散（Bordoloi et al., 2022）は MSD \(\sim t^{0.67\pm0.04}\) を与え、輸送指数 \(2/3\) を確認する。
		
		\subsection{直接的な「2/3乗則」の例}
		いくつかの研究は陽に「2/3乗則」を報告している：固着液滴の蒸発（Stauber et al., 2015）、\emph{E. coli} の表面--体積スケーリング（Kale et al., 2023）、異常拡散（Bordoloi et al., 2022）。これらは最も直接的でモデルに依存しない確認を与える。
		
		\section{統計的有意性}
		
		表1に挙げた15の独立検定をフィッシャーの方法で統合すると、30自由度で \(\chi^2 = 112.3\) となり、\(p < 10^{-20}\) に相当する。これらすべての独立したデータセットが偶然に \(2/3\)（またはその派生値）と一致する指数を示す確率は天文学的に小さい。
		
		\section{考察：若い理論に対する回顧的確認}
		
		YUCT が初めて発表されたのは 2026 年 3 月——わずか 2 か月前である。この初期段階では、すべての検証は必然的に回顧的である。これは弱点ではなく、新しい理論の通常の軌跡である。重要な科学的進歩は、新しいデータの発見ではなく、\textbf{統一}にある：これまで無関係と考えられていた数十の経験則が、すべて単一の普遍的原理 \(\varepsilon = \kappa_c \alpha K_{\mathrm{eff}}^{-2/3}\) から導かれることを示すこと。YUCT が、コーディネーションとデータ伝送を分離することで、既存のシステム（GMT、軍事、決済カード、暗号通貨）において \(K_{\mathrm{eff}} > 1\) を同定したのと同様に、今度は既存の物理的、生物的、社会則において \(\beta = 2/3\) を同定するのである。理論を確立するのは、個々のデータセットの目新しさではなく、その説明力である。
		
		\section{結論}
		
		原子スケールから宇宙論的スケールに至る十数件の独立した実験的検証が、すべて普遍的指数 \(\betaf = 2/3\) に収束する。第一原理からの理論的導出（付録 Y, L）および成功したブラインド予測（付録 S, G, Z）と合わせて、この証拠は \(\betaf = 2/3\) を自然の基本定数として、また YUCT をあらゆるスケールにわたるコーディネーションの統一枠組みとして確立する。
		
		\section*{謝辞}
		
		著者は YUCT セミナー参加者との議論、ならびに一次データソースの研究チームのオープンサイエンス実践に感謝する。本研究は Yakushev Research の支援を受けた。
		
		\section*{データの入手可能性}
		
		すべての一次データ、デジタル化手順、解析コードは YPSDC リポジトリ \url{https://github.com/Alexey-Yakushev-YUCT/YPSDC/} で入手可能である。以下に引用するテクニカルノートは YUCT のメイン DOI \url{https://doi.org/10.5281/zenodo.18444598} にアーカイブされている。
		
		\begin{thebibliography}{99}
			\bibitem{AppendixY} A. V. Yakushev, Appendix Y: Mathematical Foundations of YUCT, in \emph{YUCT} (2026).
			\bibitem{AppendixL} A. V. Yakushev, Appendix L: Fractal Coordination Error Scaling, in \emph{YUCT} (2026).
			\bibitem{AppendixC} A. V. Yakushev, Appendix C: Coordination‑Based Periodic Table, in \emph{YUCT} (2026).
			\bibitem{AppendixC2} A. V. Yakushev, Appendix C2: Universal Scaling of Phase Transitions, in \emph{YUCT} (2026).
			\bibitem{AppendixE} A. V. Yakushev, Appendix E: Coordination Genetics, in \emph{YUCT} (2026).
			\bibitem{AppendixD} A. V. Yakushev, Appendix D: Biological Predictions, in \emph{YUCT} (2026).
			\bibitem{AppendixP} A. V. Yakushev, Appendix P: Temperature Dependence of Gravity, in \emph{YUCT} (2026).
			\bibitem{AppendixM} A. V. Yakushev, Appendix M: Cosmology, in \emph{YUCT} (2026).
			\bibitem{AppendixΩ} A. V. Yakushev, Appendix Ω: Global Rotation, in \emph{YUCT} (2026).
			\bibitem{AppendixF} A. V. Yakushev, Appendix F: Coordination Economics, in \emph{YUCT} (2026).
			\bibitem{AppendixS} A. V. Yakushev, Appendix S: Negative Photon Delay, in \emph{YUCT} (2026).
			\bibitem{AppendixG} A. V. Yakushev, Appendix G: Quantum Gravity, in \emph{YUCT} (2026).
			\bibitem{AppendixZ} A. V. Yakushev, Appendix Z: Lithium Problem, in \emph{YUCT} (2026).
			\bibitem{Bordoloi2022} A. D. Bordoloi et al., \emph{Nat. Commun.} \textbf{13}, 3820 (2022).
			\bibitem{Kale2023} S. Kale et al., \emph{Phys. Biol.} \textbf{20}, 046002 (2023).
			\bibitem{Stauber2015} J. M. Stauber et al., \emph{Langmuir} \textbf{31}, 2353 (2015).
			\bibitem{Blatter1994} G. Blatter et al., \emph{Rev. Mod. Phys.} \textbf{66}, 1125 (1994).
			\bibitem{USGSNEIC} USGS NEIC Earthquake Catalog.
			\bibitem{FishBase} FishBase POPGROWTH Table.
			\bibitem{TN1} A. V. Yakushev, Technical Note 1: Anomalous Diffusion, Fragmentation, Glass Relaxation (2026).
			\bibitem{TN2} A. V. Yakushev, Technical Note 2: Cellular Surface–Volume, Crystal Growth, Superconducting Current (2026).
			\bibitem{TN3} A. V. Yakushev, Technical Note 3: Craters, Earthquakes, Droplet Evaporation (2026).
			\bibitem{TN4} A. V. Yakushev, Technical Note 4: Fish Growth, City Sizes, Metabolic Scaling (2026).
			\bibitem{TN5} A. V. Yakushev, Technical Note 5: Bacterial Scaling, Crystal Growth, Superconductivity (2026).
		\end{thebibliography}
		
	\end{multicols}
\end{document}